\documentclass[11pt,a4paper]{article}
\usepackage[T1]{fontenc}
\usepackage[utf8]{inputenc}
\usepackage[margin=24mm]{geometry}
\usepackage[hidelinks]{hyperref}
\pdfsuppressptexinfo=15
\pdfinfoomitdate=1
\hypersetup{pdftitle={Instalar o Lume em uma TV LG},pdfauthor={Projeto Lume}}
\setlength{\parindent}{0pt}
\setlength{\parskip}{5pt}
\begin{document}
\sloppy
{\LARGE\bfseries Instalar o Lume em uma TV LG}\\[4pt]
Edição pública 1.0.0 \quad|\quad PowerShell no Windows
\par\textbf{Arquivos para compartilhar.} Envie este PDF e o arquivo
\verb|com.lume.launcher_1.0.0_all.ipk|. O IPK já contém o aplicativo e o serviço auxiliar. Cada pessoa instala as ferramentas no próprio computador e faz o pareamento com a própria TV.
\section*{1. Prepare a TV e o computador}
Use uma TV LG com webOS e o aplicativo \textbf{Developer Mode} disponível na loja. O alvo inicial é webOS TV 4.x. Conecte TV e computador à mesma rede. Redes de convidados podem isolar dispositivos.

Na loja da TV, instale \textbf{Developer Mode}, entre com sua conta LG Developer, ative \textbf{Dev Mode Status} e aceite a reinicialização. Reabra o aplicativo, ative \textbf{Key Server} e anote o IPv4 da TV. Mantenha essa tela aberta durante o pareamento.

\textbf{Instale o Node.js no Windows.} Abra \url{https://nodejs.org/en/download}, escolha uma versão LTS e baixe o instalador Windows Installer (\verb|.msi|) para a arquitetura do computador (x64 para Intel ou AMD de 64 bits). Abra o arquivo, aceite a licença, mantenha selecionados o npm e a opção de adicionar ao PATH e conclua em \textbf{Install} e \textbf{Finish}.

Abra uma \textbf{nova} janela do PowerShell. Os dois comandos abaixo apenas exibem as versões instaladas. Continue quando ambos mostrarem um número de versão.
\begin{verbatim}
node -v
npm.cmd -v
\end{verbatim}
Instale agora a \href{https://webostv.developer.lge.com/develop/tools/cli-installation}{CLI oficial da LG}, confira sua versão e selecione o perfil de TV:
\begin{verbatim}
npm.cmd install -g @webos-tools/cli
ares.cmd -V
ares-config.cmd --profile tv
\end{verbatim}
Se um comando recém-instalado não for reconhecido, reabra o terminal e confira se o diretório global do npm está no PATH.
\newpage\section*{2. Cadastre e pareie a TV}
Use o cadastro interativo para evitar erros de digitação:
\begin{verbatim}
ares-setup-device.cmd
\end{verbatim}
Escolha \textbf{add}. Informe o apelido \verb|minha-tv|, o IPv4 exibido na TV, a porta \textbf{9922} e o usuário \textbf{prisoner}. Aceite os demais valores padrão e confirme o cadastro. Em seguida:
\begin{verbatim}
ares-novacom.cmd -d minha-tv --getkey
ares-device.cmd -d minha-tv --system-info
\end{verbatim}
Digite a passphrase mostrada na TV apenas no terminal. O último comando deve mostrar informações da TV. Em caso de timeout, substitua \verb|IP_DA_TV| em \verb|Test-NetConnection IP_DA_TV -Port 9922| e confira a rede. Se a porta responder e a autenticação falhar, reative o Key Server e repita a obtenção da chave.
\section*{3. Instale e abra}
Use o IPK recebido com este guia ou baixe-o da pasta \textbf{downloads} do projeto. Abra o PowerShell na pasta onde salvou o arquivo:
\begin{verbatim}
ares-install.cmd -d minha-tv .\com.lume.launcher_1.0.0_all.ipk
ares-launch.cmd -d minha-tv com.lume.launcher
\end{verbatim}
Confirme que o Lume abriu e que as setas movem a seleção. Nas próximas vezes, abra Lume pela lista de aplicativos da LG. O botão Home físico continua abrindo a interface definida pela TV.
\newpage\section*{4. Configure os recursos opcionais}
\begin{itemize}
\item \textbf{Apps e temas.} Em Demais apps, atualize a lista e aceite na TV a consulta aos aplicativos instalados. Em Preferências, escolha o tema e organize os cartões. A cidade da previsão é opcional.
\item \textbf{Câmeras.} Abra Preferências e selecione Servidores locais. Informe o endereço do Agent DVR com protocolo, IPv4 privado e porta. No campo de câmeras, informe um ou dois IDs numéricos. Abra Câmeras para ver o mosaico. O servidor precisa permitir que a TV leia os endpoints de imagem.
\item \textbf{Galeria.} No mesmo menu, informe o endereço do Immich. Ao abrir Galeria, digite na TV uma chave da sua conta com permissão para ler as fotos, álbuns e metadados desejados. Prefira HTTPS. Com HTTP, a chave trafega sem criptografia. Para encerrar o acesso, revogue a chave no Immich e limpe o endereço na TV.
\end{itemize}
\section*{5. Renove a sessão ou desinstale}
Confira o tempo restante no Developer Mode e use \textbf{Extend} enquanto a sessão estiver ativa. Ao desativar ou deixar expirar o modo de desenvolvimento, os aplicativos de desenvolvimento podem ser removidos. Para reinstalar, repita o pareamento se necessário e instale o IPK.
\begin{verbatim}
ares-install.cmd -d minha-tv --remove com.lume.launcher
ares-setup-device.cmd --remove minha-tv
\end{verbatim}
A primeira linha remove o Lume da TV. A segunda remove o cadastro da TV no computador. Para remover a CLI, use \verb|npm.cmd uninstall -g @webos-tools/cli|.

\textbf{Para reconstruir.} Na raiz do repositório, execute \verb|ares-package.cmd app service -o downloads|. Consulte também o \href{https://webostv.developer.lge.com/develop/getting-started/developer-mode-app}{guia oficial de Developer Mode da LG}.
\end{document}
